% =====================================================================
%  Chapter: Conclusion
% =====================================================================
\chapter{Conclusion}
\label{ch:conclusion}

\section{Summary of findings}
This thesis built a single, reproducible SPX pipeline --- from raw
OptionMetrics exports to a unified five-method evaluation --- and used it to
traverse the methodological spectrum of implied-volatility surface
construction.

On the parametric side, we replicated the complete toolkit of
\cite{gatheral2014arbitrage} to the digit, correcting in passing a term in
the printed inverse of Lemma~3.2, and turned it into an audited benchmark.
The benchmark quantifies the fit--arbitrage trade-off (SVI $\approx0.97$ vol
points hold-out with $6.2\%$ butterfly-violating slices; SSVI arbitrage-free
at $\approx2.8$ vol points) and localizes the difficulty: violations
concentrate in ultra-short maturities and, contrary to intuition, in the
\emph{core} moneyness region rather than the wings --- a structural
limitation of the SVI form where the ATM kink meets small total variance.
The transparent pipeline matches the proprietary OptionMetrics surface to
$0.93$ vol points at the median.

On the learning side, the derivative-constrained deep smoother --- an SSVI
prior multiplied by a positive neural corrector, trained with
exact-autodiff Durrleman and calendar penalties on a collocation grid ---
showed in controlled ablations that the prior dominates in the sparse-data
regime (RMSE $0.187$ vs $1.418$ without it) and that the constraints act as
a regularizer (beating the unconstrained variant, $0.187$ vs $0.247$) while
keeping $g\ge0$. Where per-slice SVI ceases to exist at $25\%$ quote
coverage, the surface-level smoother still reconstructs at $0.315$ vol
points. Finally, the neural operator lifts the remaining economic
constraint --- one training per day --- and its low-data transfer experiment
(real-only, synthetic-zero-shot, pretrain-plus-finetune) addresses a regime
the original paper does not: smoothing with two orders of magnitude less
data than the industrial setting.

\section{Limitations}
Three limitations bound the claims. First, several headline numbers in this
draft come from a preliminary 8-day January-2018 sample; the full 2018--2022
runs replace them before submission. Second, soft no-arbitrage constraints
\emph{reduce} violations but do not \emph{abolish} them outside the training
support \cite{chataigner2020deep}; we report residual rates rather than
claiming guarantees. Third, a small methodological asymmetry (unweighted
SSVI objective versus spread-weighted SVI) slightly disadvantages SSVI in
the current tables and is scheduled for removal.

\section{Future work}
Natural continuations include: training the full graph neural operator of
\cite{gonon2024operator} at GPU scale on the daily sample; hybrid hard--soft
constraint schemes; extending the smoother to price space with a Dupire
local-volatility extraction and half-variance regularization as in
\cite{chataigner2020deep}; uncertainty quantification via operator
ensembles; and intraday data, which would test whether the low-data
conclusions of the transfer experiment persist as the data regime approaches
the industrial one.

